% ============================================================
%  SECTION II — RELATED WORK
%  \input'd from main.tex; contains no \begin{thebibliography}.
%  Bibliography is in references.tex, \input'd at the end of main.
% ============================================================

\section{Related Work}
\label{sec:related}

\subsection{AI-Generated Image Detection Datasets}
\label{ssec:datasets}

The construction of large-scale benchmarks for distinguishing
synthetic from real imagery has a trajectory shaped by the
successive dominance of different generative architectures.

Wang~\etal~\cite{wang2020cnn} introduced \textbf{ForenSynths}, the
first widely adopted benchmark for this task, assembling images from
eleven GAN architectures to study cross-generator generalisation.
Their central finding---that a ResNet-50 trained on one GAN generator
degrades sharply when transferred to another---remains one of the
most influential results in the field and directly motivated the
cross-generator evaluation protocol we adopt here.
However, ForenSynths predates the diffusion model era and provides
no coverage of the backbone families that dominate modern synthesis.

As generative models shifted from GANs to diffusion models,
ForenSynths became inadequate.
Zhu~\etal~\cite{zhu2023genimage} addressed this with
\textbf{GenImage}, a million-scale benchmark pairing real ImageNet
images with synthetic counterparts from eight diffusion generators.
GenImage provided scale and generator diversity, but subsequent
analysis by Grommelt~\etal~\cite{grommelt2024jpeg} revealed a severe
JPEG compression disparity: real images are sourced from ImageNet
with a modal JPEG quality factor of approximately 96, whereas
generated images are stored as lossless PNG.
Detectors trained on GenImage thus partially function as JPEG
artifact classifiers rather than synthesis detectors, inflating
reported cross-generator AUROC by as much as 11~percentage
points~\cite{grommelt2024jpeg}.
Our dataset eliminates this bias class entirely by applying a single
canonical preprocessing pipeline identically to both real and AI
images.

Hong and Zhang~\cite{hong2024wildfake} introduced \textbf{WildFake},
collecting real user-generated prompts from community platforms such
as Civitai and Midjourney.
WildFake's key empirical finding---that generator \emph{architecture},
rather than fine-tuning variant, is the primary source of
generator-specific forensic bias---is directly consistent with the
UNet-vs-non-UNet generalisation gap we report and independently
validates our decision to diversify across backbone families.

The more recent \textbf{AI-GenBench}~\cite{pellegrini2025aigenbench}
adopts an incremental perspective, ordering generators by historical
release date and measuring how detection accuracy degrades as new
models arrive.
Their finding that mean detection accuracy falls from $\approx 79\%$
(2020--2021 models) to $\approx 38\%$ (2024 models) quantifies the
arms-race trajectory and directly motivates our inclusion of recent
architectures---FLUX.1-schnell, PixArt-$\Sigma$, and W\"{u}rstchen
---whose detectability profiles we report as novel benchmarks.

\subsection{Dataset Biases and the Leak Problem}
\label{ssec:bias}

The work of Grommelt~\etal~\cite{grommelt2024jpeg} is the most
direct precedent for the preprocessing design of \textbf{\datasetname}.
They identify two pervasive biases: JPEG quality factor disparity
and resolution disparity.
Because real images are typically web-sourced JPEG while generated
images are lossless PNG, a detector has access to a trivial shortcut
unrelated to synthesis artefacts.
Removing both biases produces measurable gains in cross-generator
generalisation, which Grommelt~\etal\ attribute entirely to the
suppression of shortcut learning.
Our canonical preprocessing pipeline operationalises their
recommendation at scale, applies it identically to both classes,
and verifies the result with a zero-violation audit.

Corvi~\etal~\cite{corvi2023detection} characterised the forensic
properties of diffusion model outputs in depth, demonstrating that
existing GAN-era detectors fail almost entirely on diffusion-generated
images.
Their key finding---that diffusion models introduce qualitatively
different artefacts concentrated in specific frequency bands,
compared with GAN fingerprints---informs two design decisions in
our work.
First, we use the sniff-test accuracy as a per-generator diagnostic
for artefact strength rather than as a binary pass/fail gate, since
elevated sniff scores in the absence of format asymmetry reflect
genuine synthesis fingerprints rather than preprocessing shortcuts.
Second, it grounds our discussion of why EfficientNet-B4 outperforms
ResNet-50 (Section~\ref{ssec:disc_resnet}): non-local,
frequency-distributed artefacts require larger effective receptive
fields than local $3\times 3$ convolutions provide.

The most direct motivation for our caption-grounded pairing strategy
comes from \textbf{ImageDetectBench}~\cite{jiang2024passive}.
Jiang~\etal\ observe that if AI-generated images are produced from
the same captions as their real-image counterparts, the content
distribution between the two classes becomes matched, forcing
detectors to rely on synthesis artefacts rather than semantic
shortcuts.
Their benchmark covers four dataset--generator pairings of up to
11{,}000 images each.
Our work scales this principle to 70{,}000 images across six
generators simultaneously, using captions produced by
BLIP-2~\cite{li2023blip2} from the real images themselves, and
extends the auditing framework to include canonical-format
verification and a structured sniff test.
Crucially, the identical cleaned caption is passed byte-for-byte
to all six generators, eliminating content variation between
generators by construction.

\subsection{Generative Models Used in This Work}
\label{ssec:generators}

Our dataset draws from six text-to-image architectures spanning
the main design families active in the 2022--2024 period.

\textbf{Stable Diffusion~1.5} and \textbf{SDXL}~\cite{rombach2022ldm,podell2024sdxl}
are UNet latent diffusion models operating in the compressed
latent space of a pretrained variational autoencoder, with the
denoising network conditioned on CLIP text embeddings.
SDXL extends SD~1.5 with a three-times-larger UNet, a second text
encoder, and native 1024$\times$1024 resolution.
Both remain the most widely deployed open-weight generators and
serve as the UNet-family anchor in our cross-generator experiments.

\textbf{PixArt-$\Sigma$}~\cite{chen2024pixart} replaces the UNet
backbone with a Diffusion Transformer (DiT) architecture, enabling
efficient generation at up to 4K resolution.
The shift from convolutional to attention-based denoising produces
qualitatively different artefact distributions, reflected in its
high cross-generator detectability (AUROC~$= 0.993$ when held out).

\textbf{FLUX.1-schnell}~\cite{flux2024bfl} is a rectified flow
Transformer operating natively in latent space, representing the
most recent generation of open-weight models as of mid-2024.
Its four-step distilled inference regime produces fewer iterative
sampling artefacts than earlier diffusion models, placing it in
an intermediate detectability position relative to other non-UNet
generators.

\textbf{Kandinsky~2.2}~\cite{razzhigaev2023kandinsky} employs a
two-stage prior--decoder architecture: a prior network maps text
to image embeddings in CLIP space, and a UNet latent diffusion
decoder synthesises the image.
The separation of semantic and visual decoding produces strongly
distinctive artefacts; Kandinsky achieves the second-highest
sniff-test score (0.958) and near-perfect cross-generator
detectability (AUROC~$= 0.993$) when held out.

\textbf{W\"{u}rstchen}~\cite{pernias2023wuerstchen} employs a
hierarchical two-stage latent compression scheme with an extreme
42$\times$ spatial compression ratio, distinguishing it from
all other generators in the dataset.
Its artefact profile is the most challenging for the sniff test
(0.843, the only healthy-zone score) yet achieves near-perfect
cross-generator detectability (AUROC~$= 0.992$), a combination
analysed in detail in Section~\ref{ssec:disc_wuerstchen}.

\subsection{Detection Methods}
\label{ssec:detectors}

The standard detection pipeline, established by Wang~\etal's
CNNSpot~\cite{wang2020cnn}, fine-tunes a pretrained ResNet-50 as a
binary real-vs-fake classifier.
Despite its simplicity, CNNSpot-style classifiers remain the
dominant evaluation baseline in dataset papers, and we follow this
convention with three representative architectures.

Ojha~\etal~\cite{ojha2023universal} showed that the CNNSpot approach
fails catastrophically when trained on GAN images and evaluated on
diffusion images.
Their proposed fix---frozen CLIP features with nearest-neighbour
classification---substantially improves cross-family generalisation
and motivates our inclusion of ViT-B/16 (which shares CLIP's
attention-based pretraining) among our three baselines.

Corvi~\etal~\cite{corvi2023detection} demonstrated that frequency-domain
analysis of diffusion model outputs reveals artefacts at specific
spectral bands, motivating detection approaches based on DCT
coefficient statistics or wavelet decompositions.
These approaches are complementary to the spatial-domain classifiers
evaluated here and represent a natural avenue for future work on
\textbf{\datasetname}.

Jiang~\etal~\cite{jiang2024passive} argue that passive detectors are
fundamentally limited by the artefact distribution present in their
training data, and that watermark-based proactive detection
consistently outperforms passive approaches under post-processing
perturbations.
Because passive shortcuts are provably absent from \textbf{\datasetname},
it constitutes a fair evaluation environment for both passive and
proactive methods---an important property not shared by benchmarks
that carry compression disparities.

\subsection{Positioning of This Work}
\label{ssec:position}

Table~\ref{tab:dataset_comparison} summarises how \textbf{\datasetname}
compares against the five most closely related prior benchmarks.
The key distinctions are threefold.
First, we are the only dataset to apply a \emph{single canonical
preprocessing function}---identical in every detail for real and AI
images---that provably eliminates both JPEG quality and resolution
asymmetry simultaneously, verified by a zero-violation
canonical-format audit.
Second, our caption-grounded pairing ensures that all six generators
receive byte-identical prompts derived from the same real image,
eliminating semantic content as a confounding variable while our
pair-level splitting ensures that no prompt appears in more than one
split.
Third, we are the first paired benchmark to span all three principal
backbone families---UNet LDM, Diffusion Transformer, and hierarchical
two-stage---enabling the cross-family generalisation measurement
that reveals the UNet-vs-non-UNet gap.

\begin{table*}[t]
\centering
\caption{Comparison of AI-generated image detection datasets.
``Paired'' denotes that real and AI images share a common content
anchor (caption or class label).
``Leak-free'' denotes that an identical preprocessing pipeline was
applied to both classes.
``Multi-family'' denotes coverage of more than one backbone
architecture family (UNet, Transformer, hierarchical).}
\label{tab:dataset_comparison}
\setlength{\tabcolsep}{8pt}
\begin{tabular}{@{}lcccccl@{}}
\toprule
Dataset & Images & Generators & Paired & Leak-free & Multi-family & Venue \\
\midrule
ForenSynths~\cite{wang2020cnn}              & 362K & 11 (GANs) & $\times$     & $\times$     & $\times$     & CVPR 2020 \\
GenImage~\cite{zhu2023genimage}             & 1.3M & 8         & $\times$     & $\times$     & $\times$     & NeurIPS 2023 \\
WildFake~\cite{hong2024wildfake}            & 84K  & 17+       & $\times$     & $\times$     & $\checkmark$ & arXiv 2024 \\
AI-GenBench~\cite{pellegrini2025aigenbench} & 310K & 20+       & $\times$     & $\times$     & $\checkmark$ & IJCNN 2025 \\
ImageDetectBench~\cite{jiang2024passive}    & 44K  & 4         & $\checkmark$ & $\times$     & $\times$     & arXiv 2024 \\
\midrule
\textbf{Ours (\datasetname)} & \textbf{70K} & \textbf{6} & $\checkmark$ & $\checkmark$ & $\checkmark$ & -- \\
\bottomrule
\end{tabular}
\end{table*}
